(1) 给定 \(\mathbf M=M_z(r)\hat z\)，体束缚电流为 \(\mathbf J_b=\nabla\times\mathbf M=-dM_z/dr\,\hat{\boldsymbol{\phi}}=2M_0\alpha r/R^2\,\hat{\boldsymbol{\phi}}\)。外侧表面束缚电流为 \(\mathbf K_b=\mathbf M\times\hat{\mathbf r}=M_0(1-\alpha)\hat{\boldsymbol{\phi}}\)，注意 \(\hat z\times\hat r=\hat\phi\)。这些电流在周向流动而非沿柱轴流动；其符号由空间梯度和交叉乘积决定。

(2) 由无自由电流的 Maxwell 方程 \(\nabla\times\mathbf H=0\)、\(\nabla\cdot\mathbf B=0\)，利用轴向平移和旋转对称性，可取 \(\mathbf H=0\)，满足无穷远磁场衰减条件、无磁标势外源及跨界面切向 \(H_z\) 连续条件。于是柱内 \(\mathbf B(r)=\mu_0\mathbf M(r)=\mu_0M_0(1-\alpha r^2/R^2)\hat z\)，柱外 \(\mathbf B=0\)。边界内外切向 \(B_z\) 的跳跃由面束缚电流给出，不能误以为必须连续。

(3) 将各个半径上的体周向电流按半径积分，相当于穿过定方位的径向半平面，得到每单位轴长上的体周向片电流 \(K_{\rm vol}=\int_0^R J_{b,\phi}(r)\,dr=M_0\alpha\)。表面贡献 \(K_{\rm surf}=M_0(1-\alpha)\)，两者相加为 \(M_0\)。若从总磁化最大值的圆心处向外积分，所得结果对应维持 \(B_z(0)=\mu_0M_0\) 的总同向环流；在这里体与表面电流应相加而非相互抵消。反之若 \(\alpha=0\)，仅表面束缚电流流动。

(4) 在固定磁化分布且不追究建立磁化的物质内部能量时，可计算真空形式的磁感应场量 \(\mathcal U_B=\int B^2/(2\mu_0)\,dA\)（每单位轴长）。柱外为零，柱内积分得 \(\mathcal U_B=\pi\mu_0M_0^2R^2\int_0^1u(1-\alpha u^2)^2\,du=\pi\mu_0M_0^2R^2[1/2-\alpha/2+\alpha^2/6]\)。此量是场的 \(B^2/(2\mu_0)\) 积分，不应不加说明地称为具有固定磁化永久磁体的全部热力学自由能。令 \(\alpha\to0\)，退化为 \(B=\mu_0M_0\) 的均匀内部场与 \(\mathcal U_B=\pi R^2\mu_0M_0^2/2\)。
